% REPLACES: online_appendix.tex lines 3038-3083 ("\subsection{C.6c. Exact reserve events ..." to "... searches found none.")
% Numbers in this fragment: \NUM{reserve_event_sampled_weak_max}. The value 6.280 in
% the current text came from the earlier sweep and sat next to the old endpoint
% p_L = h - 1/m. Under the zero cost that endpoint is h, and revenue \rho p/2 rises
% in p on (r,h) wherever the high cost stays out, so the sampled weak maximum must
% come from the new sweep. The formulas at any one such reserve do not change.
\subsection{C.6c. Exact reserve events and admissible-bid
outcomes}\label{oa-c-reserve-events}

This exercise runs the C.6 node solver, with the continuation identity
and event handling above, on the event grid alone. It leaves the full
sweep untouched and exists so that the event candidates and the outcome
measures can be checked without the hour-long scan.

\textbf{Inputs.} Both economies at both strengths, on the event grid of
C.6: \(0\), \(\ell\), \(r\), \(h\), the band edges in the class economy,
the declared reserve alternatives, the endpoint \(p_L=h\) at every strength and \(p_H\)
where \(\ell<p_H<r\), each with the three one-sided offsets on each
side that stays in the domain, and the reserve
\NUM{reserve_event_sampled_weak_max} used as the sampled weak maximum
in the current sweep.

\textbf{Method and acceptance.} At \(p_L=h\) in the weak economy, the row
must show the equality-event classification of the zero-cost comparison at
the domain endpoint, a single
accepted full-order continuation whose status names the equality \(B(\mu)=0\) and the
tie rule rather than a strict inequality, \(\mathsf E=\rho\), sale
probability \(\rho/2\), \(\mathsf C_2=0\), and
\(\mathcal R_T=\rho h/2\) from both the row and a fifty-digit
evaluation; the three offsets below the event must classify the comparison as
strictly positive. No offset above \(h\) belongs to the domain, and no row
classifies that side. At \(p_H\)
in the strong economy, the row must show the equality-event
classification of the ceiling with gross profit \(mc_H/M>0\) at the lowest
belief and a
strictly positive full-order bound, \(\tau=M\), \(x^*=1\),
\(\alpha_H=1/2\), \(\alpha_L=e^{-2/b}/2\), and entry and revenue
agreeing with the fifty-digit evaluation; where \(p_H\) falls outside
its defining regime the node is an ordinary reserve node and must not be
classified as an event. At the sampled reserve
\NUM{reserve_event_sampled_weak_max}, when that reserve lies in \((r,h)\) with
the high cost excluded,
\(\mathsf E=\rho\), sale \(\rho/2\), \(\mathsf C_2=0\), and
seller revenue \(\rho\) times that reserve, divided by two: preparation there does not mean two
admissible bidders, because the incumbent is excluded. On every row the
union identity for \(\mathsf S\) and the ordering
\(0\le\mathsf C_2\le\mathsf A\le\mathsf E\le1\) hold, and the payoff
oracle error is within tolerance. Deduplication follows the C.6 rule,
and every count in the manifest comes from the ledger.

\textbf{Outputs.} \texttt{numerics/reserve\_events.csv}, with the
continuation schema of C.6 and an added \texttt{event\_classification}
text column, and \texttt{numerics/reserve\_event\_ranges.csv} with the
range schema of C.6. Feed: the exact-event paragraphs of the paper
appendix and the online summary of C.6. Event candidates improve sampled
maxima; they prove neither global optimality nor envelope completeness,
and a node without an accepted candidate is a node where the declared
searches found none.
